\chapter{CƠ SỞ LÝ THUYẾT}

\section{Phương pháp nghiên cứu}
Các phương pháp được sử dụng trong quá trình thực hiện đề tài bao gồm:

\begin{itemize}
\item
	\textbf{Phân tích yêu cầu:} Thu thập và phân tích quy trình đặt lịch
	thực tế tại các Barbershop.
\item
	\textbf{Mô hình hóa hệ thống:} Sử dụng sơ đồ Use Case, sơ đồ tuần tự
	(Sequence Diagram) và sơ đồ thực thể liên kết (ERD) để thiết kế hệ
	thống.
\item
	\textbf{Phát triển hệ thống:}

	\begin{itemize}
	\item
		Frontend: HTML5, CSS3, Vanilla JavaScript (sử dụng Fetch API).
	\item
		Backend: PHP Native (kiến trúc RESTful API).
	\end{itemize}
\item
	\textbf{Kiểm thử chức năng:} Thực hiện kiểm thử độc lập các API bằng
	Postman trước khi tích hợp vào giao diện.
\end{itemize}

\section{Công nghệ được sử dụng}
\subsection{Frontend}
\textbf{HTML, CSS và JavaScript} là bộ ba ngôn ngữ nền tảng cốt lõi để
xây dựng giao diện cho mọi trang web. Trong đó, HTML đóng vai trò xây
dựng cấu trúc và bố cục trang; CSS đảm nhận việc định dạng, trang trí
và làm đẹp giao diện; còn JavaScript chịu trách nhiệm xử lý các sự kiện,
tạo hiệu ứng và tương tác với người dùng.

\begin{itemize}
\item
	\textbf{Ưu điểm nổi bật:}

	\begin{itemize}
	\item
		Khởi chạy trực tiếp trên mọi trình duyệt web mà không cần cài đặt
		thêm môi trường phức tạp.
	\item
		Tốc độ tải trang nhanh, nhẹ do không bị phụ thuộc vào các thư viện
		hay framework cồng kềnh của bên thứ ba.
	\item
		Dễ dàng gỡ lỗi (debug) và chỉnh sửa trực tiếp thông qua công cụ của
		trình duyệt (Developer Tools).
	\end{itemize}
\item
	\textbf{Lý do sử dụng trong đề tài:}

	\begin{itemize}
	\item
		Phù hợp để xây dựng giao diện website Barbershop nhanh chóng, thân
		thiện và đáp ứng tốt trên các thiết bị.
	\item
		Giúp người lập trình nắm vững bản chất cốt lõi của việc thao tác với
		phần tử web (DOM) và các sự kiện của người dùng.
	\end{itemize}
\end{itemize}

\textbf{Fetch API} là một giao diện lập trình ứng dụng (API) hiện đại
được tích hợp sẵn trong JavaScript, cho phép trình duyệt gửi các yêu
cầu mạng (HTTP Request) đến máy chủ để lấy hoặc gửi dữ liệu.

\begin{itemize}
\item
	\textbf{Ưu điểm nổi bật:}

	\begin{itemize}
	\item
		Xử lý dữ liệu theo cơ chế bất đồng bộ (Asynchronous), giúp trang web
		không bị đơ hoặc gián đoạn khi đang chờ dữ liệu từ máy chủ.
	\item
		Cú pháp viết code ngắn gọn, rõ ràng và dễ đọc hơn rất nhiều so với
		các công nghệ cũ (như XMLHttpRequest).
	\end{itemize}
\item
	\textbf{Lý do sử dụng trong đề tài:}

	\begin{itemize}
	\item
		Được sử dụng làm cầu nối để giao diện web (Frontend) gửi dữ liệu đặt
		lịch hoặc từ khóa tìm kiếm lên máy chủ (Backend).
	\item
		Giúp trang web cập nhật dữ liệu mới (ví dụ: hiển thị danh sách giờ
		trống, danh sách khách hàng) ngay lập tức mà không cần phải tải lại
		(reload) toàn bộ trang web.
	\end{itemize}
\end{itemize}

\subsection{Backend}
\textbf{PHP} là một ngôn ngữ lập trình kịch bản chạy ở phía máy chủ
(Server-side), được thiết kế chuyên biệt để phát triển các ứng dụng web
động.

\begin{itemize}
\item
	\textbf{Ưu điểm nổi bật:}

	\begin{itemize}
	\item
		Dễ học, dễ triển khai và có khả năng tích hợp cực kỳ tốt với cơ sở
		dữ liệu MySQL.
	\item
		Xử lý linh hoạt các luồng dữ liệu do người dùng gửi lên từ các biểu
		mẫu (Form).
	\item
		Cộng đồng hỗ trợ lớn, hoạt động cực kỳ ổn định trên các máy chủ ảo
		cục bộ như XAMPP.
	\end{itemize}
\item
	\textbf{Lý do sử dụng trong đề tài:}

	\begin{itemize}
	\item
		Đóng vai trò là bộ não xử lý nghiệp vụ của hệ thống (ví dụ: tính
		toán giờ trống, kiểm tra trùng lịch, xác thực thông tin đăng nhập).
	\item
		Tiếp nhận yêu cầu từ JavaScript, truy vấn cơ sở dữ liệu và trả về
		kết quả dưới dạng chuỗi dữ liệu (JSON) để giao diện web hiển thị.
	\end{itemize}
\end{itemize}

\subsection{Cơ sở dữ liệu}
\textbf{MySQL} là một trong những hệ quản trị cơ sở dữ liệu quan hệ
(RDBMS) mã nguồn mở phổ biến nhất thế giới. Dữ liệu trong MySQL được
lưu trữ dưới dạng các bảng (Table) có cấu trúc chặt chẽ và liên kết với
nhau thông qua các khóa (Key).

\begin{itemize}
\item
	\textbf{Ưu điểm nổi bật:}

	\begin{itemize}
	\item
		Quản lý dữ liệu có cấu trúc một cách an toàn, chính xác và nhất
		quán.
	\item
		Tốc độ truy vấn dữ liệu nhanh, hỗ trợ tốt các câu lệnh tìm kiếm, lọc
		dữ liệu phức tạp thông qua ngôn ngữ SQL.
	\item
		Dễ dàng quản lý, sao lưu và phục hồi dữ liệu thông qua công cụ trực
		quan phpMyAdmin.
	\end{itemize}
\item
	\textbf{Lý do sử dụng trong đề tài:}

	\begin{itemize}
	\item
		Rất phù hợp với đặc thù dữ liệu của hệ thống đặt lịch, nơi cần sự
		liên kết chặt chẽ giữa bảng Khách hàng, bảng Dịch vụ, bảng Thợ cắt
		tóc và bảng Lịch hẹn.
	\item
		Đảm bảo tính toàn vẹn của dữ liệu, không để xảy ra tình trạng mất
		mát hoặc sai lệch thông tin lịch hẹn của khách hàng.
	\end{itemize}
\end{itemize}

\section{Các công cụ hỗ trợ}
\textbf{Visual Studio Code:} Phần mềm soạn thảo mã nguồn.

\textbf{XAMPP:} Môi trường máy chủ ảo cục bộ (Localhost) để chạy mã
nguồn PHP và lưu trữ cơ sở dữ liệu.

\textbf{Postman:} Công cụ hỗ trợ kiểm thử luồng truyền nhận dữ liệu.

\textbf{phpMyAdmin:} Công cụ quản trị cơ sở dữ liệu trực quan trên nền
tảng web.

\section{Tổng kết chương}
Chương 2 đã trình bày phương pháp thực hiện, các công nghệ cốt lõi và
công cụ hỗ trợ được sử dụng để xây dựng hệ thống đặt lịch. Đây là cơ sở
để chuyển sang phần phân tích, thiết kế và triển khai kỹ thuật ở các
chương tiếp theo.
